\section{Conclusion}

In this study, we introduce \textsc{MIRAGE-Bench} for systematically identifying and evaluating hallucinations in interactive LLM-based agent scenarios. \textsc{MIRAGE-Bench} categorizes hallucinations based on their unfaithfulness to task instructions, interaction history, or environment observations. An extensive audit across various environments reveals that even advanced proprietary models continue to face significant faithfulness issues when it comes to realistic risk settings that are common in the wild.
Our findings underscore that hallucination remains a persistent and widespread challenge, with proprietary models showing only modest improvements over open-source counterparts. This suggests that beyond simple scaling or instruction tuning, targeted alignment techniques specific to interactive settings are crucial.
Given the static nature of our evaluation framework, future research should consider dynamic, rollout-based assessments to better reflect real-world interactions. Expanding the benchmark to embodied agents, multi-modal contexts, and specialized tasks will further clarify hallucination phenomena. As agents become increasingly complex, continuous refinement of evaluation frameworks will be essential.

